\subsection{Kết luận chương 1}

Chương 1 đã trình bày một cách hệ thống các công nghệ và cơ sở lý thuyết nền tảng phục vụ cho việc xây dựng hệ thống bỏ phiếu điện tử dựa trên công nghệ chuỗi khối được đề xuất trong đồ án. Cụ thể, mục 1.1 giới thiệu tổng quan về Blockchain với các khái niệm cốt lõi, đặc điểm, cấu trúc dữ liệu, phân loại, cơ chế đồng thuận và hợp đồng thông minh. Các mục 1.3, 1.4, 1.5 và 1.6 lần lượt trình bày các cơ sở mật mã học cần thiết: hàm băm và chữ ký số, mật mã đường cong elliptic dựa trên bài toán ECDLP, lược đồ chữ ký EC-Schnorr và lược đồ chữ ký mù do David Chaum đề xuất. Trong đó, hai khối kiến thức trọng tâm nhất, đóng vai trò là hai trụ cột công nghệ của đồ án, được trình bày chi tiết trong mục 1.2 và mục 1.7.

\textit{Hyperledger Fabric} (mục 1.2) được lựa chọn làm nền tảng blockchain cho hệ thống bỏ phiếu điện tử bởi đây là nền tảng blockchain consortium mã nguồn mở do Linux Foundation phát triển, có kiến trúc mô-đun linh hoạt với các thành phần chính gồm Peer, Orderer, Certificate Authority, Channel và Ledger. Đặc biệt, Hyperledger Fabric cung cấp ba ưu điểm phù hợp với yêu cầu của hệ thống bầu cử: \textit{tính riêng tư và phân quyền} thông qua cơ chế Membership Service Provider và Channel cho phép giới hạn quyền truy cập của từng nhóm thành viên; \textit{hiệu năng giao dịch cao} nhờ kiến trúc Execute-Order-Validate tách biệt khâu thực thi với khâu sắp xếp; và \textit{tính lập trình mạnh mẽ} thông qua Chaincode hỗ trợ nhiều ngôn ngữ lập trình phổ thông. Những đặc tính này đáp ứng đồng thời các yêu cầu về tính minh bạch, toàn vẹn, hiệu năng và phân quyền truy cập của một hệ thống bỏ phiếu điện tử quy mô lớn.

\textit{Chữ ký mù tập thể dựa trên EC-Schnorr} (mục 1.7) là công cụ mật mã cốt lõi của đồ án, được xây dựng từ sự kết hợp giữa lược đồ chữ ký mù (mục 1.6) và lược đồ chữ ký EC-Schnorr (mục 1.5). Lược đồ này cho phép ban tổ chức bầu cử (gồm nhiều thành viên đóng vai trò người ký tập thể) chứng thực phiếu bầu của cử tri mà không biết nội dung lá phiếu, đồng thời không thể truy lại mối liên kết giữa cử tri và lá phiếu sau khi kết quả được công bố. Lược đồ đáp ứng đầy đủ ba thuộc tính an toàn cốt lõi: \textit{tính mù} bảo vệ nội dung lá phiếu khỏi ban tổ chức, \textit{tính không thể giả mạo} chống lại các phiếu bầu gian lận đã được chứng minh chặt chẽ trong mô hình ROM thông qua quy giảm về bài toán ECDLP, và \textit{tính không thể liên kết} đảm bảo nặc danh cho cử tri ngay cả khi tất cả thành viên ban tổ chức thoả hiệp lại với nhau. Đặc biệt, kích thước chữ ký tập thể cố định không phụ thuộc số lượng người ký và quá trình xác minh chỉ cần một phép kiểm tra duy nhất bằng khoá công khai tập thể -- đây là những đặc tính có ý nghĩa quyết định đối với một hệ thống triển khai trên blockchain, nơi mỗi byte dữ liệu và mỗi phép tính đều phát sinh chi phí giao dịch trực tiếp.

Sự kết hợp giữa nền tảng Hyperledger Fabric và lược đồ chữ ký mù tập thể dựa trên EC-Schnorr chính là cơ sở công nghệ then chốt cho phép giải quyết đồng thời ba yêu cầu vốn được xem là đối nghịch nhau trong các hệ thống bỏ phiếu điện tử truyền thống: tính minh bạch -- tính riêng tư -- hiệu năng cao. Các kết quả lý thuyết được trình bày trong chương 1 sẽ được vận dụng trực tiếp để thiết kế kiến trúc hệ thống, mô hình dữ liệu và các giao thức nghiệp vụ chi tiết của hệ thống bỏ phiếu điện tử đề xuất trong các chương tiếp theo của đồ án.
